Penampang horizontal dicirikan oleh lenyapnya pemetaan vertikalnya. Dalam situasi ini, menurut
contoh tersebut,
turunan vertikal sama dengan

\mavergleichskettedisp
{\vergleichskette
{ \nabla_\partial (f)
}
{ =} { f' + f { \frac{  g' }{  2g } }
}
{ } {
}
{ } {
}
{ } {}
}
{}{}{.}
Pertama-tama,

\mavergleichskette
{\vergleichskette
{ f
}
{ = }{ 0
}
{  }{
}
{    }{
}
{   }{
}
}
{}{}{}
merupakan suatu solusi. Sekarang misalkan

\mavergleichskette
{\vergleichskette
{ f(t_0)
}
{ \neq }{ 0
}
{  }{
}
{    }{
}
{   }{
}
}
{}{}{}
dan karena itu sifat ini juga berlaku pada suatu lingkungan terbuka dari $t_0$. Pada lingkungan tersebut,

\mavergleichskettedisp
{\vergleichskette
{  { \frac{  f' }{ f } }
}
{ =} { { \frac{  g' }{  2g  } }
}
{ } {
}
{ } {
}
{ } {
}
}
{}{}{.}
Metode penyelesaian persamaan diferensial linear menghasilkan

\mavergleichskettedisp
{\vergleichskette
{   \ln  \betrag {  f  }
}
{ =} { c+ { \frac{   1  }{  2 } }  \ln   g
}
{ } {
}
{ } {
}
{ } {
}
}
{}{}{}
dan dengan demikian

\mavergleichskettedisp
{\vergleichskette
{ \betrag {  f  }
}
{ =} { e^c e^{ { \frac{   1  }{  2 } }  \ln  g }
}
{ =} { e^c  \sqrt{g}
}
{ } {
}
{ } {
}
}
{}{}{}
atau

\mavergleichskettedisp
{\vergleichskette
{ f
}
{ =} { d \sqrt{g}
}
{ } {
}
{ } {
}
{ } {
}
}
{}{}{}
\zusatzklammer {dengan

\mavergleichskette
{\vergleichskette
{ d
}
{ \in }{  \R
}
{  }{
}
{    }{
}
{   }{
}
}
{}{}{}} {} {,}
yang juga dapat diperiksa secara langsung.
